% =====================================================================
% Reemplaza el bloque
%   \section*{Energía  potencial del campo gravitacional} ... \end{align}
% de main.tex (la deducción con E_pot = F·r).
% Requiere: amsmath, amssymb, siunitx, booktabs (ya cargados en main.tex).
% =====================================================================

\subsection{Energía potencial del campo gravitacional}
\label{subsec:energia_potencial_gravitacional}

\paragraph{Punto de partida: trabajo y energía potencial.}
Recordemos que una fuerza \emph{constante} $\boldsymbol{F}$ que actúa sobre un
cuerpo mientras este experimenta un desplazamiento $\Delta\boldsymbol{r}$
realiza un trabajo igual al producto escalar de los dos vectores,
\begin{equation}
  W = \boldsymbol{F}\cdot\Delta\boldsymbol{r}
    = \lvert\boldsymbol{F}\rvert\,\lvert\Delta\boldsymbol{r}\rvert\cos\theta \,,
  \label{eq:trabajo_producto_escalar}
\end{equation}
donde $\theta$ es el ángulo entre $\boldsymbol{F}$ y $\Delta\boldsymbol{r}$.
Cerca de la superficie, la gravedad hace sobre un cuerpo que sube una altura
$h$ el trabajo $W = -mgh$, y por eso definimos la energía potencial $U = mgh$: el
cambio de energía potencial es el negativo del trabajo que hace la gravedad,
\begin{equation}
  \Delta U = -\,W_{\text{grav}} \,.
  \label{eq:def_cambio_epot}
\end{equation}
Usaremos la misma definición para el campo gravitacional a cualquier
distancia, pero ahora con la referencia en el infinito:
$U(\infty) = 0$.

\paragraph{La dificultad: la fuerza no es constante.}
Sobre una masa $m$ colocada en la posición $\boldsymbol{r}$, medida desde el
centro de la Tierra (masa $M$), actúa la fuerza gravitacional
\begin{equation}
  \boldsymbol{F}(\boldsymbol{r})
  = -\,\frac{GMm}{r^{3}}\,\boldsymbol{r}
  = -\,\frac{GMm}{r^{2}}\,\hat{\boldsymbol{r}} \,,
  \label{eq:fuerza_grav_epot}
\end{equation}
cuya magnitud, $GMm/r^{2}$, cambia con la distancia. La
ecuación~\eqref{eq:trabajo_producto_escalar} no se puede aplicar de una sola
vez a todo el recorrido. En particular, \emph{no} es correcto escribir
$U = \boldsymbol{F}\cdot\boldsymbol{r}$: el vector
$\boldsymbol{r}$ no es un desplazamiento del proceso y la fuerza no es la
misma a lo largo del camino. Lo que haremos es dividir el camino en tramos
cortos y acotar el trabajo en cada uno.

\paragraph{Paso 1: el trabajo en un tramo radial.}
Llevemos la masa $m$ en línea recta, alejándola del centro, desde la distancia
$r_{a}$ hasta la distancia $r_{b} > r_{a}$. Dividimos el camino en $n$ tramos
iguales de longitud
\begin{equation}
  \Delta r = \frac{r_{b} - r_{a}}{n} \,,
  \label{eq:ancho_tramo}
\end{equation}
con extremos
\begin{equation}
  r_{0} = r_{a}, \quad
  r_{i} = r_{a} + i\,\Delta r \quad (i = 0, 1, \dots, n), \quad
  r_{n} = r_{b} \,.
  \label{eq:extremos_tramos}
\end{equation}
En el tramo $i$, que va de $r_{i-1}$ a $r_{i}$, el desplazamiento es
$\Delta\boldsymbol{r}_{i} = \Delta r\,\hat{\boldsymbol{r}}$, dirigido hacia
afuera, mientras que la fuerza gravitacional apunta hacia el centro. El ángulo
entre ellos es $\theta = 180^{\circ}$ y $\cos\theta = -1$. Si en ese tramo la
fuerza tuviera el valor constante que tiene a una distancia $r$, el
trabajo~\eqref{eq:trabajo_producto_escalar} sería
\begin{equation}
  \boldsymbol{F}\cdot\Delta\boldsymbol{r}_{i}
  = \left(-\,\frac{GMm}{r^{2}}\,\hat{\boldsymbol{r}}\right)
    \cdot\left(\Delta r\,\hat{\boldsymbol{r}}\right)
  = -\,\frac{GMm}{r^{2}}\,\Delta r\,
    \left(\hat{\boldsymbol{r}}\cdot\hat{\boldsymbol{r}}\right)
  = -\,\frac{GMm}{r^{2}}\,\Delta r \,,
  \label{eq:trabajo_tramo_fuerza_fija}
\end{equation}
porque $\hat{\boldsymbol{r}}\cdot\hat{\boldsymbol{r}} = 1$.

\paragraph{Paso 2: cotas del trabajo en cada tramo.}
Dentro del tramo $i$ la distancia está entre $r_{i-1}$ y $r_{i}$, de modo que
la magnitud de la fuerza disminuye desde $GMm/r_{i-1}^{2}$, al comienzo,
hasta $GMm/r_{i}^{2}$, al final. Por lo tanto, el trabajo verdadero $W_{i}$
de la gravedad en ese tramo, que es negativo, tiene una magnitud mayor que la
que tendría con la fuerza más débil del tramo y menor que la que tendría con
la más intensa:
\begin{equation}
  \frac{GMm}{r_{i}^{2}}\,\Delta r
  \;<\; -\,W_{i} \;<\;
  \frac{GMm}{r_{i-1}^{2}}\,\Delta r \,.
  \label{eq:cotas_trabajo_tramo}
\end{equation}
Geométricamente, $-W_{i}$ es el área bajo la curva $F(r) = GMm/r^{2}$ entre
$r_{i-1}$ y $r_{i}$, y las dos cotas son las áreas de los rectángulos de base
$\Delta r$ que quedan por debajo y por encima de la curva.

\paragraph{Paso 3: una resta de fracciones con las mismas cotas.}
Restando fracciones,
\begin{equation}
  \frac{1}{r_{i-1}} - \frac{1}{r_{i}}
  = \frac{r_{i} - r_{i-1}}{r_{i-1}\,r_{i}}
  = \frac{\Delta r}{r_{i-1}\,r_{i}} \,.
  \label{eq:resta_fracciones}
\end{equation}
Como $0 < r_{i-1} < r_{i}$, al multiplicar esta desigualdad por $r_{i-1}$ y
por $r_{i}$ se obtiene $r_{i-1}^{2} < r_{i-1}\,r_{i} < r_{i}^{2}$; tomando
inversos, que invierten el orden de los números positivos, y multiplicando
por $\Delta r > 0$,
\begin{equation}
  \frac{\Delta r}{r_{i}^{2}}
  \;<\; \frac{1}{r_{i-1}} - \frac{1}{r_{i}} \;<\;
  \frac{\Delta r}{r_{i-1}^{2}} \,.
  \label{eq:cotas_resta}
\end{equation}
La diferencia $1/r_{i-1} - 1/r_{i}$ queda entre las mismas cotas que
$-W_{i}/(GMm)$ en~\eqref{eq:cotas_trabajo_tramo}.

\paragraph{Paso 4: sumar todos los tramos.}
El trabajo total de la gravedad es la suma de los trabajos de los tramos,
$W_{\text{grav}} = W_{1} + W_{2} + \dots + W_{n}$. Definimos las sumas
inferior y superior
\begin{equation}
  S_{\text{inf}} = GMm\,\Delta r \sum_{i=1}^{n} \frac{1}{r_{i}^{2}} \,,
  \qquad
  S_{\text{sup}} = GMm\,\Delta r \sum_{i=1}^{n} \frac{1}{r_{i-1}^{2}} \,.
  \label{eq:sumas_inf_sup}
\end{equation}
Sumando las desigualdades~\eqref{eq:cotas_trabajo_tramo} para
$i = 1, \dots, n$,
\begin{equation}
  S_{\text{inf}} \;<\; -\,W_{\text{grav}} \;<\; S_{\text{sup}} \,.
  \label{eq:cotas_trabajo_total}
\end{equation}
Sumando las desigualdades~\eqref{eq:cotas_resta}, multiplicadas por $GMm$,
el término del medio se cancela en cadena (suma \emph{telescópica}):
\begin{align}
  \sum_{i=1}^{n}\left(\frac{1}{r_{i-1}} - \frac{1}{r_{i}}\right)
  &= \left(\frac{1}{r_{0}} - \frac{1}{r_{1}}\right)
   + \left(\frac{1}{r_{1}} - \frac{1}{r_{2}}\right)
   + \dots
   + \left(\frac{1}{r_{n-1}} - \frac{1}{r_{n}}\right) \nonumber\\
  &= \frac{1}{r_{0}} - \frac{1}{r_{n}}
   = \frac{1}{r_{a}} - \frac{1}{r_{b}} \,,
  \label{eq:suma_telescopica}
\end{align}
porque cada $1/r_{i}$ intermedio aparece una vez sumando y otra restando. Por
lo tanto,
\begin{equation}
  S_{\text{inf}} \;<\; GMm\left(\frac{1}{r_{a}} - \frac{1}{r_{b}}\right)
  \;<\; S_{\text{sup}} \,.
  \label{eq:cotas_valor_exacto}
\end{equation}

\paragraph{Paso 5: las cotas se juntan.}
La diferencia entre las dos sumas también es telescópica:
\begin{align}
  S_{\text{sup}} - S_{\text{inf}}
  &= GMm\,\Delta r \sum_{i=1}^{n}
     \left(\frac{1}{r_{i-1}^{2}} - \frac{1}{r_{i}^{2}}\right)
   = GMm\,\Delta r \left(\frac{1}{r_{0}^{2}} - \frac{1}{r_{n}^{2}}\right)
   \nonumber\\
  &= GMm\,\frac{r_{b} - r_{a}}{n}
     \left(\frac{1}{r_{a}^{2}} - \frac{1}{r_{b}^{2}}\right) .
  \label{eq:diferencia_cotas}
\end{align}
Esta diferencia es inversamente proporcional al número de tramos: al
multiplicar $n$ por $10$ se divide por $10$, y se puede hacer tan pequeña
como se quiera. Según~\eqref{eq:cotas_trabajo_total}
y~\eqref{eq:cotas_valor_exacto}, los números $-W_{\text{grav}}$ y
$GMm\,(1/r_{a} - 1/r_{b})$, que no dependen de $n$, están \emph{siempre}
dentro del mismo intervalo $(S_{\text{inf}},\, S_{\text{sup}})$, cuyo ancho
tiende a cero. La única posibilidad es que sean iguales:
\begin{equation}
  W_{\text{grav}} = -\,GMm\left(\frac{1}{r_{a}} - \frac{1}{r_{b}}\right) .
  \label{eq:trabajo_grav_r1_r2}
\end{equation}
El signo menos indica que, al alejarse la masa, la gravedad se opone al
movimiento. En la tabla~\ref{tab:cotas_trabajo} se ven las cotas para
$r_{a} = R$ y $r_{b} = 4R$, en unidades de $GMm/R$; allí el valor exacto es
$1 - \tfrac{1}{4} = 0{,}75$.

\begin{table}[htbp]
  \centering
  \caption{Sumas inferior y superior del trabajo, en unidades de $GMm/R$, para
    llevar la masa de $r_{a} = R$ a $r_{b} = 4R$. El valor exacto,
    $0{,}7500$, queda siempre entre las dos, y la diferencia se divide por
    $10$ cada vez que $n$ se multiplica por $10$.}
  \label{tab:cotas_trabajo}
  \begin{tabular}{@{}rccc@{}}
    \toprule
    $n$ & $S_{\text{inf}}$ & $S_{\text{sup}}$ & $S_{\text{sup}} - S_{\text{inf}}$ \\
    \midrule
    4   & $0{,}4828$ & $1{,}1859$ & $0{,}7031$ \\
    40  & $0{,}7158$ & $0{,}7861$ & $0{,}0703$ \\
    400 & $0{,}7465$ & $0{,}7535$ & $0{,}0070$ \\
    \bottomrule
  \end{tabular}
\end{table}

\paragraph{Paso 6: otros caminos.}
Si el desplazamiento de un tramo es perpendicular a $\hat{\boldsymbol{r}}$,
es decir, se mueve sobre una esfera centrada en $M$, entonces $\theta =
90^{\circ}$ y, por~\eqref{eq:trabajo_producto_escalar}, la gravedad no
realiza trabajo. Cualquier camino de $r_{a}$ a $r_{b}$ puede aproximarse por
escalones radiales y escalones sobre esferas; solo los radiales contribuyen,
y su suma es la misma que la del camino recto. El
resultado~\eqref{eq:trabajo_grav_r1_r2} depende únicamente de las distancias
inicial y final, y por eso tiene sentido hablar de una energía potencial que
depende solo de $r$.

\paragraph{Paso 7: la energía potencial.}
Con la definición~\eqref{eq:def_cambio_epot},
\begin{equation}
  U(r_{b}) - U(r_{a})
  = -\,W_{\text{grav}}
  = GMm\left(\frac{1}{r_{a}} - \frac{1}{r_{b}}\right) .
  \label{eq:cambio_epot_r1_r2}
\end{equation}
Ahora alejamos la masa indefinidamente: $r_{b} \to \infty$, de modo que
$1/r_{b} \to 0$, y escogemos $U(\infty) = 0$. Llamando $r$ a la
distancia inicial $r_{a}$,
\begin{equation}
  0 - U(r) = \frac{GMm}{r}
  \quad\Longrightarrow\quad
  \boxed{U(r) = -\,\frac{GMm}{r}} \,.
  \label{eq:Epot_grav}
\end{equation}
La energía potencial es negativa a cualquier distancia finita: la masa está en
un \emph{pozo} de energía, y para sacarla del campo hay que darle al menos la
energía $GMm/r$.

\paragraph{Comprobación: cerca de la superficie se recupera $mgh$.}
Con $r_{a} = R_{T}$ y $r_{b} = R_{T} + h$ en~\eqref{eq:cambio_epot_r1_r2},
\begin{equation}
  U(R_{T} + h) - U(R_{T})
  = GMm\,\frac{(R_{T} + h) - R_{T}}{R_{T}\,(R_{T} + h)}
  = \frac{GM}{R_{T}^{2}}\,m\,h\,\frac{1}{1 + h/R_{T}} \,.
  \label{eq:epot_cerca_exacta}
\end{equation}
Para $h \ll R_{T}$, $1 + h/R_{T} \approx 1$, y con
$g = GM/R_{T}^{2} \approx \SI{9.8}{\meter\per\second\squared}$,
\begin{equation}
  U(R_{T} + h) - U(R_{T}) \approx m\,g\,h \,,
  \label{eq:epot_cerca_mgh}
\end{equation}
que es la energía potencial que usamos cerca de la superficie. La
expresión~\eqref{eq:Epot_grav} la contiene como caso particular.
